\chapter*{阅读指引}

分类器在一批训练邮件上答对很多次，仍可能只记住了这批输入；即使当前测试错误很少，换一批
训练邮件也可能得到另一种规则。理论部分要说明，有限样本在什么条件下足以支持未来风险的
判断。为此，我们同时规定观测怎样产生、允许哪些决策、如何计损失、算法能看到什么，以及
成功要相对于哪个基准、在哪些随机变化中成立。

基础三章已经区分任务、模型、损失及反馈，数学两卷则提供组合维度、集中不等式、随机复杂度、
信息量、分布距离和优化条件。本部分把这些工具接到具体学习问题：候选空间怎样受控，
样本怎样影响选解，经验评价怎样接近总体风险。涉及网络时，只建立推导所需的形式与条件；
结构设计、反向传播和训练实践由卷五展开。

《机器学习理论基础》以邮件二分类建立风险、PAC标准和归纳偏置，说明拟合经验为何还缺少
未来保证。随后先从有限候选扩展到无限假设空间，再从二分类扩展到多分类，分别核对统计
样本量、停机及计算资源。《一般学习的可学习性》把零一损失换为一般损失，从全类一致估计
进一步转向算法稳定性。《深度学习的
可学习性》考察权重大小、初始化附近的梯度、更新所得参数及随机参数分布，分别在什么条件下
支持未来风险的分析。《迁移学习的可学习性》
最后改变训练与评价分布，先判断目标预测能否由已有信息确定，再分析覆盖、共享机制和环境
抽样如何支持有限样本保证。

六章复用共同的风险、概率量词和可测性约定，也逐步改变标签、损失、选解方式及数据来源。
每次复用结论，都要重新核对这些对象和条件。所得保证约束的是允许问题中的风险或超额风险，
还没有判断损失是否覆盖实际后果、比较基准是否可接受，也不能自动覆盖允许范围之外的环境。
下一部分将这些使用问题逐项纳入评价与处置。
